% Compiler avec XeLaTeX (Overleaf : Menu > Compilateur > XeLaTeX),
% ou avec : tectonic 10-reflexes-oral-kholle.tex
\documentclass[11pt,a4paper]{article}
\usepackage[margin=17mm,top=16mm,bottom=18mm]{geometry}
\usepackage{fontspec}
\setmainfont{texgyrepagella}[Extension=.otf,UprightFont=*-regular,
  BoldFont=*-bold,ItalicFont=*-italic,BoldItalicFont=*-bolditalic]
\setsansfont{texgyreheros}[Extension=.otf,UprightFont=*-regular,
  BoldFont=*-bold,ItalicFont=*-italic,BoldItalicFont=*-bolditalic]
\usepackage[french]{babel}
\usepackage{amsmath,amssymb,xcolor,tabularx,fancyhdr}
\usepackage[colorlinks=true,urlcolor=teal,linkcolor=teal]{hyperref}
\definecolor{ink}{HTML}{153342}
\definecolor{teal}{HTML}{147D80}
\definecolor{muted}{HTML}{55646D}
\definecolor{wash}{HTML}{EEF5F5}
\definecolor{line}{HTML}{D6E4E6}
\hypersetup{pdftitle={10 réflexes pour conduire son oral de concours},
  pdfsubject={Fiche pratique - khôlles et oraux scientifiques de prépa},
  pdfauthor={Projet X-IA}}
\setlength{\parindent}{0pt}
\setlength{\parskip}{0pt}
\setlength{\emergencystretch}{2em}
\pagestyle{fancy}
\fancyhf{}
\renewcommand{\headrulewidth}{0pt}
\fancyfoot[L]{\sffamily\footnotesize\color{muted}X-IA\enspace /\enspace Fiche d'entraînement à l'oral}
\fancyfoot[R]{\sffamily\footnotesize\color{muted}\thepage\,/\,2}
\newcommand{\repere}[4]{%
  \begin{minipage}{\linewidth}
  \begin{tabularx}{\linewidth}{@{}p{10mm}@{\hspace{3mm}}X@{}}
  {\sffamily\bfseries\fontsize{20}{22}\selectfont\color{teal}#1}&
  {\sffamily\bfseries\large\color{ink}#2}\par\vspace{2mm}
  #3\par\vspace{1mm}
  {\small\color{teal}\textit{#4}}
  \end{tabularx}
  \end{minipage}\par\vspace{4.3mm}}
\newcommand{\bandeau}[2]{%
  \begingroup\setlength{\fboxsep}{4mm}%
  \colorbox{wash}{\parbox{\dimexpr\linewidth-8mm\relax}{%
  {\sffamily\bfseries\color{ink}#1}\par\vspace{2mm}#2}}%
  \endgroup}
\begin{document}
\color{ink}
{\sffamily\bfseries\small\color{teal}PRÉPA SCIENTIFIQUE\enspace /\enspace KHÔLLES ET CONCOURS}\par
\vspace{3mm}
{\sffamily\bfseries\fontsize{27}{30}\selectfont 10 réflexes pour\par conduire son oral}\par
\vspace{3mm}
{\large Faire comprendre son idée, l'ancrer dans le cours,\par puis la défendre et l'ajuster dans le dialogue.}\par
\vspace{3mm}
{\small\color{muted}Des repères pour progresser, sans garantie de note : le fond scientifique reste essentiel.}
\par\vspace{5mm}

\repere{01}{Reformuler la question}
{Précise les objets, les données et ce qu'il faut établir. En physique, fixe le système, le référentiel et les hypothèses du modèle.}
{« Je cherche à montrer que… ; l'hypothèse qui semble utile est… »}

\repere{02}{Annoncer l'idée avant le calcul}
{Choisis une piste et explique son intérêt. Donne un plan court ; une liste de méthodes possibles sans choix ne fait pas avancer.}
{« Je vais majorer ce terme pour me ramener à une série connue. »}

\repere{03}{Rattacher explicitement l'idée au cours}
{Nomme le résultat pertinent et précise ce qu'il permet de conclure. Si une définition ou une preuve est demandée, donne-la avec précision.}
{« Le théorème spectral me donne ici une base orthonormée de vecteurs propres. »}

\repere{04}{Vérifier les hypothèses}
{Fais le lien entre les conditions du théorème et les données. Vérifie aussi que les objets manipulés existent et que les opérations sont licites.}
{« La matrice est réelle et symétrique : le théorème spectral s'applique. »}

\repere{05}{Rendre le raisonnement visible}
{Explique les étapes décisives, distingue conjecture et preuve, garde des notations stables. Organise le tableau et laisse le jury voir ce que tu écris.}
{« Ce calcul suggère une limite ; il reste maintenant à prouver la convergence. »}

\vfill
\bandeau{LE FIL À GARDER}{%
\centering\sffamily\bfseries\large
Objectif $\longrightarrow$ Idée $\longrightarrow$ Cours $\longrightarrow$ Hypothèses $\longrightarrow$ Preuve}

\newpage
{\sffamily\bfseries\small\color{teal}PENDANT L'ÉCHANGE}\par
\vspace{3mm}
{\sffamily\bfseries\fontsize{23}{26}\selectfont Écouter, justifier, rebondir}\par
\vspace{4.5mm}

\repere{06}{Accueillir une relance sans paniquer}
{Écoute la question jusqu'au bout. Une demande de justification ne signifie pas que tout est faux : réponds au point précis, puis reprends ton raisonnement.}
{« Vous me demandez pourquoi cette majoration est valable : voici l'argument. »}

\repere{07}{Transformer l'indication en action}
{Reformule l'aide reçue, puis utilise-la concrètement. Être réactif, c'est identifier la prochaine étape utile, avec le temps de réflexion nécessaire.}
{« Votre indication me conduit à étudier les sommes partielles ; je les note $S_n$. »}

\repere{08}{Corriger une erreur et garder le cap}
{Repère la ligne fautive, explique la correction et ses conséquences. Si ton argument reste valable, justifie-le calmement ; suis une réorientation demandée.}
{« J'ai divisé par une quantité qui peut être nulle. Je distingue les deux cas. »}

\repere{09}{Gérer un blocage et son temps}
{Dis précisément ce qui manque. Reviens à la définition, teste un cas simple ou exploite une hypothèse inutilisée. Évite le silence prolongé et les calculs sans objectif.}
{« Il me manque une borne indépendante de $n$ ; je cherche une majoration uniforme. »}

\repere{10}{Conclure et contrôler le résultat}
{Réponds explicitement à la question. Vérifie un cas limite ; en physique, contrôle unités, signe et ordre de grandeur. Distingue ce qui est établi de ce qui reste ouvert.}
{« J'ai démontré la convergence ; la valeur de la somme reste à déterminer. »}

\bandeau{EXEMPLE ORIGINAL : RÉAGIR À UNE OBJECTION}{%
\small
\textbf{Sujet.} Étudier la série de terme général $u_n=1/[n(n+1)]$, pour $n\geq1$.\par
\textbf{Élève.} « $u_n\to 0$, donc $\sum u_n$ converge. »\par
\textbf{Colleur.} « Que dit cette règle pour $u_n=1/n$ ? »\par
\textbf{Élève.} « La série harmonique diverge : mon critère était insuffisant.
Je reprends $u_n=\dfrac{1}{n(n+1)}$, pour $n\geq1$ :
$u_n=\dfrac1n-\dfrac1{n+1}$. Ainsi,
$S_N=1-\dfrac1{N+1}\to1$. Par la définition de la convergence d'une série,
la série converge et sa somme vaut $1$. »\par\vspace{1mm}
\textbf{À observer :} l'élève utilise le contre-exemple, corrige son idée,
mobilise une définition du cours et termine la preuve.}

\vfill
{\footnotesize\color{muted}
\textbf{Appuis et usage.} Fiche pédagogique originale ; les phrases et le dialogue sont des exemples d'entraînement.
Repères de jury :
\href{https://concoursminesponts.fr/wp-content/uploads/2025/10/Rapport-Oral-2025-VF.pdf#page=7}{Mines-Ponts, rapport oral 2025, p. 3 et 6--8}
(méthode, hypothèses, interaction) ;
\href{https://www.concours-centrale-supelec.fr/sites/default/files/documents/Rapports_Jury_2025_MP.pdf#page=78}{Centrale-Supélec, rapport MP 2025, p. 78--80}
(cours, clarté, écoute et tableau). Les modalités dépendent de l'épreuve.}
\end{document}
